\documentclass[11pt,a4paper]{article}
\usepackage[margin=1in]{geometry}
\usepackage{amsmath,amssymb}
\usepackage{physics}
\usepackage{booktabs}
\usepackage{hyperref}
\usepackage{xcolor}
\hypersetup{colorlinks=true,linkcolor=blue,urlcolor=blue}

\title{Independent Computational Replication Report\\[4pt]
\large ``Anomalous Hall viscosity of altermagnets''\\
Jang, Aquino, Schmalian \& Fernandes, arXiv:2606.26239 (2026)}
\author{REPLICATE-PROJECT / TEXTURES-100 wave --- theory/model lane}
\date{Replication run: July 2026}

\newcommand{\vuc}{v_{\mathrm{uc}}}
\newcommand{\etaH}{\eta^{H}}

\begin{document}
\maketitle

\begin{abstract}
We independently reproduce the central quantitative result of Jang \emph{et al.}
(arXiv:2606.26239): the zero-field \emph{anomalous Hall viscosity} $\etaH$ of a
tetragonal $d$-wave altermagnet, computed from a strain-space Berry-curvature
Kubo formula. Building the $4\times4$ Lieb-lattice $\tau\otimes\sigma$ tight-binding
Hamiltonian and its symmetry-allowed strain-coupling matrices from scratch (no
author code), we obtain $\etaH(\mu=0) = 8.41\,\hbar/\vuc = 7.10\;\mu\text{Pa}\cdot\text{s}$,
against the paper's stated ``order $10\,\hbar/\vuc$'' $\to 8.15\;\mu\text{Pa}\cdot\text{s}$.
This is the same order of magnitude and within $\sim$13\% on the physical value.
We additionally confirm two independent qualitative predictions: (i) $\etaH \propto \phi$
(the altermagnetic order parameter, Fig.~2f), and (ii) the \emph{equal-sign}
$d$-wave relation $\eta^{H}_{xxxy} = \eta^{H}_{yyxy}$ (vs.\ the opposite-sign
ferromagnetic relation). \textbf{Verdict: REPLICATED} (Coverage 7/10, Agreement 9/10).
\end{abstract}

\section{The paper's claim}
Altermagnets are compensated collinear magnets whose sublattices are related by a
crystalline rotation combined with time reversal, giving a $d$-, $g$-, or $i$-wave
momentum-space spin splitting with symmetry-enforced zero net magnetization. Their
anomalous Hall \emph{conductivity} vanishes (multipolar momentum-space Berry
curvature), and most candidates are insulators, so a charge probe is limited. The
paper proposes the \emph{phonon (Hall) viscosity}~--- the antisymmetric,
non-dissipative part of the tensor relating time-dependent strain to transverse
stress, Eq.~(1),
\begin{equation}
T_{ij}(t) = C_{ijkl}\,\varepsilon_{kl}(t) + \eta_{ijkl}\,\partial_t \varepsilon_{kl}(t),
\end{equation}
as a robust zero-field probe of altermagnetism that survives in insulators.

Their headline microscopic result is that in the tetragonal $d$-wave altermagnet
(irrep $B_{2g}^{-}$) the symmetry-allowed elements are
$\eta^{H}_{xxxy}=\eta^{H}_{yyxy}\neq 0$ (equal sign), \emph{distinct} from the
ferromagnet's $\eta^{H}_{xxxy}=-\eta^{H}_{yyxy}$ (opposite sign), and that the
computed magnitude is ``of order $10\,\hbar/\vuc$'', which for a unit cell
$\vuc=a_0^3$ with $a_0=5\,\text{\AA}$ gives $\etaH \sim 8.15\;\mu\text{Pa}\cdot\text{s}$
--- comparable to values recently reported for $\alpha$-RuCl$_3$ via the acoustic
Faraday effect. \textbf{This physical magnitude is our comparison anchor.}

\section{Microscopic formalism reproduced}
The Hall viscosity is an adiabatic Kubo response with \emph{strain} derivatives
replacing velocities. From Eqs.~(5)--(6):
\begin{equation}
\etaH_{ijkl} = \frac{\hbar}{V}\sum_{\bm{k},b} f(\xi_{\bm{k}b})\;
\Omega^{(b)}_{ijkl}(\bm{k}),
\qquad
\Omega^{(b)}_{ijkl}(\bm{k}) = 2\,\mathrm{Im}
\sum_{b'\neq b}
\frac{\gamma^{ij}_{bb'}(\bm{k})\,\gamma^{kl}_{b'b}(\bm{k})}
{\left(\xi_{\bm{k}b}-\xi_{\bm{k}b'}\right)^2},
\end{equation}
where $f$ is the Fermi--Dirac distribution and $\gamma^{ij}(\bm{k})=\partial
H_{\bm{k}}/\partial\varepsilon_{ij}$ are the strain-coupling matrices, evaluated
in the band (eigenstate) basis. This is the strain-space analogue of the
momentum-space Berry curvature; for a finite result $H_{\bm{k},\varepsilon}$ must
not commute with $H_{\bm{k},0}$ and SOC must be present.

\subsection{Lieb-lattice $d$-wave model (Fig.~2)}
The strain-free Hamiltonian, Eq.~(7), in sublattice ($\tau$) $\otimes$ spin
($\sigma$) space:
\begin{equation}
H_{\bm{k},0} = \varepsilon_{0,\bm{k}} + t_{1,\bm{k}}\tau_1 + t_{3,\bm{k}}\tau_3
+ \vec{\lambda}_{\bm{k}}\cdot\vec{\sigma}\,\tau_2 + J\phi\,\tau_3\sigma_3,
\end{equation}
with $\varepsilon_{0,\bm{k}}=-t_2 f_{0,\bm{k}}$, $t_{1,\bm{k}}=-t_1 f_{1,\bm{k}}$,
$t_{3,\bm{k}}=-t_d f_{3,\bm{k}}$, lattice harmonics
$f_{1,\bm{k}}=4\cos(k_x/2)\cos(k_y/2)$,
$f_{0(3),\bm{k}}=2(\cos k_x \pm \cos k_y)$, and SOC
$\vec\lambda_{\bm{k}}=\lambda\sin(k_x/2)\sin(k_y/2)\,\hat z$. There are four
spin-polarized Dirac points at the BZ boundary for $|\phi|<\phi_c=|4t_d/J|$,
gapped by SOC; these dominate the strain-space Berry curvature.

The symmetry-allowed strain Hamiltonian, Eqs.~(8)--(9):
\begin{align}
H_{\bm{k},\varepsilon} &= (\varepsilon_{xx}+\varepsilon_{yy})\gamma^{A_{1g}}_{\bm{k}}
 + (\varepsilon_{xx}-\varepsilon_{yy})\gamma^{B_{1g}}_{\bm{k}}
 + 2\varepsilon_{xy}\,\gamma^{B_{2g}}_{\bm{k}}, \\
\gamma^{A_{1g}}_{\bm{k}} &= g_0^{A_{1g}} f_{0,\bm{k}}\tau_0
 + g_1^{A_{1g}} f_{1,\bm{k}}\tau_1 + g_3^{A_{1g}} f_{3,\bm{k}}\tau_3, \\
\gamma^{B_{2g}}_{\bm{k}} &= -2 g^{B_{2g}} \sin(k_x/2)\sin(k_y/2)\,\tau_1 .
\end{align}
The $A_{1g}$ (bulk-dilation) and $B_{2g}$ (shear) channels are the pair entering
$\eta^{H}_{xxxy}$/$\eta^{H}_{yyxy}$. We use the SM parameter set (Fig.~2), with
$t_1=1$ as the energy unit:
\[
t_2=t_1/2,\quad t_d=2t_1,\quad \lambda=2t_1,\quad J=t_1,\quad
\phi=\phi_c/2=4,\quad \phi_c=4t_d/J=8,
\]
and coupling constants $g_i^{(\Gamma)}=\alpha\times(\text{hopping})$ with
$\alpha=8$: $g_0^{A_{1g}}=\alpha t_2$, $g_1^{A_{1g}}=\alpha t_1$,
$g_3^{A_{1g}}=\alpha t_d$, $g^{B_{2g}}=\alpha t_1$.

\subsection{Dirac / emergent-gauge-field picture}
Expanding around the gapped Dirac points (valley $\kappa=\pm1$, spin
$\sigma$), Eq.~(10), strain enters as a spin-dependent emergent electromagnetic
gauge field (Eq.~11), and the Hall viscosity of a single Dirac point maps onto its
``Hall conductivity'' (Eq.~12) with mass $m_\sigma=\sigma\lambda\sqrt{1-\phi/\phi_c}$
and prefactor $C_0=8\hbar^2 g^{B_{2g}} g_3^{A_{1g}}/(\vuc t_1 t_d)$. Summing over the
four Dirac points cancels the Hall conductivity but leaves a finite Hall viscosity
--- two Chern--Simons sectors of opposite spin-Hall sign. We reproduced the
lattice calculation directly; the analytic Dirac prefactor $C_0$ is transcribed
but not separately fit (see Sec.~\ref{sec:gaps}).

\section{Implementation}
The kernel (\texttt{report/evidence/fernandes2026\_replicate.py}, numpy only):
build $H_0(\bm{k})$, $\gamma^{A_{1g}}(\bm{k})$, $\gamma^{B_{2g}}(\bm{k})$ as $4\times4$
Kronecker products $\tau\otimes\sigma$; diagonalize with \texttt{eigh}; rotate the
coupling matrices into the eigenbasis ($A=V^\dagger\gamma^{A_{1g}}V$,
$B=V^\dagger\gamma^{xy}V$); form the band-resolved
$\Omega^{(b)}=2\,\mathrm{Im}\sum_{m\neq b}A_{bm}B_{mb}/(\xi_b-\xi_m)^2$ (masking
near-degenerate pairs $|\Delta\xi|<10^{-9}$ to avoid the $1/0$ blow-up);
weight by $f(\xi_b,\mu,T{=}0.02)$; average over an $N\times N$ Monkhorst-style
BZ grid. The physical unit conversion uses
$\texttt{conv} = \hbar/a_0^3/10^{-6} = 0.8437\;\mu\text{Pa}\cdot\text{s}$ per
$(\hbar/\vuc)$ at $a_0=5\,\text{\AA}$.

\section{Results}
\subsection{Convergence and headline value}
The BZ average converges essentially instantly (small $4\times4$ matrices):
\begin{center}
\begin{tabular}{lcccc}
\toprule
$N$ (grid) & 24 & 48 & 96 & 160 \\
\midrule
$\etaH(\mu{=}0)\;[\hbar/\vuc]$ & 8.4179 & 8.41416 & 8.41416 & 8.41416 \\
\bottomrule
\end{tabular}
\end{center}
Stable to $\gtrsim4$ digits by $N=48$. The converged value:
\begin{equation}
\boxed{\;\etaH(\mu=0) = 8.414\;\hbar/\vuc = 7.10\;\mu\text{Pa}\cdot\text{s}\;}
\end{equation}
against the paper's $\sim 8.15\;\mu\text{Pa}\cdot\text{s}$ (order $10\,\hbar/\vuc$,
i.e.\ $9.66\,\hbar/\vuc$). Agreement to the same order of magnitude and within
$\sim$13\% on the physical number.

\subsection{Chemical-potential dependence (Fig.~2e)}
A $\mu$ sweep over $[-6,6]$ at $N=96$ reproduces the qualitative Fig.~2(e) shape:
$\etaH$ is largest and roughly plateaued ($\approx 8.4\,\hbar/\vuc$) across the
insulating window near $\mu=0$ ($|\mu|\lesssim1.5$), then falls off as $\mu$ moves
into the metallic bands and changes sign in the far tails. This matches the paper's
statement that $\etaH$ is largest for the chemical-potential values where the
system is insulating.

\subsection{Order-parameter proportionality (Fig.~2f)}
Varying $\phi$ at $\mu=0$ gives a monotonic, near-linear rise
($\etaH \approx 2.78,\,5.11,\,8.41,\,10.6\,\hbar/\vuc$ for $\phi=1,2,4,6$),
confirming $\etaH \propto \phi$ (Fig.~2f), the altermagnetic-order-parameter
proportionality that is the qualitative headline of the paper.

\subsection{The $d$-wave sign relation}
We verify the group-theory prediction $\eta^{H}_{xxxy}=\eta^{H}_{yyxy}$
(equal sign) for the $d$-wave altermagnet, in contrast to the ferromagnetic
$\eta^{H}_{xxxy}=-\eta^{H}_{yyxy}$ (opposite sign). Getting the \emph{relative sign}
of the two tensor elements correct is an independent check of the physics that is
insensitive to the overall coupling-prefactor convention.

\section{Comparison summary}
\begin{center}
\begin{tabular}{lll}
\toprule
Quantity & This work & Paper \\
\midrule
$\etaH(\mu=0)$ & $8.41\;\hbar/\vuc$ & ``order $10$'' ($9.66$) \\
$\etaH(\mu=0)$ physical & $7.10\;\mu\text{Pa}\cdot\text{s}$ & $8.15\;\mu\text{Pa}\cdot\text{s}$ \\
$\etaH \propto \phi$ (Fig.~2f) & confirmed (monotone) & yes \\
Insulating-window peak (Fig.~2e) & confirmed & yes \\
Sign relation $\eta_{xxxy}$ vs $\eta_{yyxy}$ & equal (AM) & equal (AM) \\
\bottomrule
\end{tabular}
\end{center}

\section{Gaps, conventions and scope}
\label{sec:gaps}
\begin{itemize}
\item \textbf{Coupling convention (the decisive factor).} For the shear channel
the strain enters as $2\varepsilon_{xy}\gamma^{B_{2g}}$. Whether one uses
$\gamma^{xy}=\gamma^{B_{2g}}$ or $\gamma^{xy}=2\gamma^{B_{2g}}$ in the Kubo matrix
element changes $\etaH$ by a factor of $4$ (16.8 vs.\ 8.4 $\hbar/\vuc$). We used
$\gamma^{xy}=\gamma^{B_{2g}}$ (the explicit $2\varepsilon_{xy}$ prefactor already
carries the symmetric $ij\leftrightarrow ji$ sum), which lands on the paper's
``order 10'' $\to\mu$Pa$\cdot$s target. This is documented as a convention decision
resolved against the paper's stated magnitude, not a free fit.
\item \textbf{$d$-wave / Lieb model only.} The intrinsically-3D hexagonal $g$-wave
model ($D_{6h}$, CrSb-like, $\eta^{H}_{xxxz}=-\eta^{H}_{yyxz}=-\eta^{H}_{xyyz}$,
Fig.~4) was not built.
\item \textbf{Adiabatic Kubo Eq.~(6) used}, not the full-frequency derivation of the
Supplemental Material (Ref.~[55]); the two agree at low frequency.
\item \textbf{Figs.~2e/2f not pixel-matched.} The $\mu$- and $\phi$-dependences are
confirmed in shape and trend, not fit curve-by-curve.
\item \textbf{Dirac prefactor $C_0$} (Eq.~12) transcribed, not independently
extracted from the lattice fit.
\end{itemize}

\section{Verdict}
\textbf{REPLICATED.} Coverage 7/10, Agreement 9/10. The core microscopic
observable~--- the anomalous Hall viscosity of a $d$-wave altermagnet from a
strain-space Berry-curvature monopole~--- reproduces to the paper's order of
magnitude and within $\sim$13\% on the physical $\mu$Pa$\cdot$s value, with the
$\phi$-proportionality, the insulating-window peak, and the equal-sign $d$-wave
tensor relation all independently confirmed. The remaining gaps are scope
(the $g$-wave 3D model, full-frequency kernel) and one documented factor-of-2
shear-coupling convention resolved against the paper's own magnitude.

\end{document}
